% TOP_PROBLEM: {"problemId": "problem.dixmier-unitarisability-problem", "canonicalTitle": "Dixmier unitarisability problem for discrete groups", "releaseRank": 376, "primaryDomain": "algebra_representation_category", "primaryDomainLabel": "Algebra, representation & category theory", "displayStatus": "open_with_solved_subcases", "releaseStatus": "open_with_solved_subcases"}
% !TeX program = lualatex
% ATTEMPT_STATUS: unresolved
\documentclass[11pt]{article}
\usepackage[margin=25mm]{geometry}
\usepackage{fontspec}
\setmainfont{Latin Modern Roman}
\usepackage{amsmath,amssymb,mathtools}
\usepackage{unicode-math}
\setmathfont{Latin Modern Math}
\usepackage{xurl,hyperref}
\hypersetup{colorlinks=true,urlcolor=blue}
\setcounter{secnumdepth}{0}
\setlength{\parindent}{0pt}
\setlength{\parskip}{5pt}
\setlength{\emergencystretch}{3em}
\begin{document}
\section*{\texorpdfstring{Dixmier unitarisability problem for discrete groups}{Dixmier unitarisability problem for discrete groups}}\label{376}
\subsection{Definitions and mathematical statement}
A discrete group $G$ is unitarisable if every homomorphism $\pi:G\to\operatorname{GL}(\mathcal H)$ on every complex Hilbert space, with $\sup_g\|\pi(g)\|<\infty$, has a bounded invertible operator $S$ such that
\[
 (S\pi(g)S^{-1})^*(S\pi(g)S^{-1})=I\qquad(g\in G).
\]
Amenability means existence of a positive normalized left-invariant mean on $\ell^\infty(G)$. The conjectured implication is that unitarisability forces amenability; equivalently, every nonamenable discrete group has at least one uniformly bounded representation not similar to a unitary representation. No countability assumption on $G$, separability assumption on $\mathcal H$, or prescribed similarity bound on $S$ is included. Boundedness of each individual $\pi(g)$ is weaker than the required uniform bound over the group.
\par\smallskip\textit{Formulation and concept references:} \href{https://arxiv.org/pdf/0902.4585v1}{[S1]}.

\subsection{Short English statement}
If every uniformly bounded Hilbert-space representation of a discrete group can be made unitary by a change of coordinates, must the group be amenable?
\subsection{Research attempt}
\paragraph{Scope and source check (27 September 2026).}
The difficult implication is from unitarisability to amenability. The
opposite implication has a short averaging proof, reconstructed below
with the direction of invariance and the similarity bound explicit.
The quantifiers include arbitrary groups and Hilbert spaces; a theorem
only about finite-dimensional representations would be much weaker.

Monod--Ozawa [S1] construct nonunitarisable representations after passing
from a nonamenable group to certain abelian extensions and deduce
nonunitarisability for certain Burnside groups. Their primary abstract
and text do not assert that every original nonamenable group already
has such a representation. Passing to an extension is precisely a scope
change that cannot be suppressed. This attempt supplies reductions and
complete subcases, not the missing general construction.

We use inner products linear in their first argument. For a uniformly
bounded representation put $M=\sup_g\|\pi(g)\|$, so $M\ge1$ and
$\|\pi(g)^{-1}\|\le M$ as well. The latter follows from the presence
of $g^{-1}$ in the same group and is crucial for a nondegenerate average.

\paragraph{An invariant positive operator is the exact certificate.}
\textbf{Lemma 376.1.} The representation $\pi$ is unitarisable if and
only if there is a bounded positive operator $Q$ and constants
$0<a\le b<\infty$ with
\[
 aI\le Q\le bI,
 \qquad \pi(g)^*Q\pi(g)=Q\quad(g\in G).
 \tag{376.1}
\]
In that case $S=Q^{1/2}$ is a unitariser with
$\|S\|\|S^{-1}\|\le\sqrt{b/a}$.

\emph{Proof.} If $S\pi(g)S^{-1}$ is unitary, multiplication of its
unitarity equation on the left by $S^*$ and on the right by $S$ gives
(376.1) with $Q=S^*S$. Conversely take the positive square root of
$Q$. Direct multiplication gives
\[
 (S\pi(g)S^{-1})^*(S\pi(g)S^{-1})
   =S^{-1}\pi(g)^*Q\pi(g)S^{-1}=I.
\]
The spectral bounds give the asserted norm estimate.\hfill$\square$

Positivity without a lower bound is not sufficient. For example a
positive operator with eigenvalues tending to zero can define an
invariant degenerate-in-the-limit norm whose change of coordinates has
unbounded inverse. The problem asks for an equivalent Hilbert norm,
not merely an invariant sesquilinear form.

\paragraph{A complete proof that amenability implies unitarisability.}
From the specified left-invariant mean obtain a right-invariant mean
$m_R$ by composing functions with inversion: if
$\check f(g)=f(g^{-1})$, set $m_R(f)=m_L(\check f)$.
A right translate $f(gh)$ becomes the left translate
$\check f(h^{-1}g)$, so this really is right invariant.

Define a sesquilinear form by
\[
 [\xi,\eta]=m_R\bigl(g\mapsto
                \langle\pi(g)\xi,\pi(g)\eta\rangle\bigr).
\]
The mean extends complex linearly from real bounded functions. Positivity
and the inequalities
\[
 M^{-2}\|\xi\|^2\le\|\pi(g)\xi\|^2\le M^2\|\xi\|^2
\]
give
$M^{-2}\|\xi\|^2\le[\xi,\xi]\le M^2\|\xi\|^2$.
The polarized form is bounded, and the Hilbert-space representation
theorem for bounded forms produces a positive $Q$ with
$[\xi,\eta]=\langle Q\xi,\eta\rangle$. Right invariance gives
\[
 [\pi(h)\xi,\pi(h)\eta]
 =m_R\bigl(g\mapsto\langle\pi(gh)\xi,\pi(gh)\eta\rangle\bigr)
 =[\xi,\eta].
\]
Thus (376.1) holds with $a=M^{-2}$ and $b=M^2$, and the similarity
constant is at most $M^2$.

No countable sum, invariant probability measure on $G$, or separability
assumption was used. A finitely additive mean suffices. Replacing it
by an ordinary uniform average over an infinite group would be undefined.
This argument proves the easy direction at the full cardinality scope
of the catalogue, but assumes the amenability one would need to deduce
in the other direction.

\paragraph{Why finite-dimensional representations cannot witness the obstruction.}
Every uniformly bounded representation into $\operatorname{GL}_d(\mathbb C)$
is unitarisable, for \emph{every} group. Indeed, its image and its
inverses are bounded in a finite-dimensional matrix space. Its closure
$K$ is compact. The uniform inverse bound prevents singular matrices in
that closure: $\|A\xi\|\ge M^{-1}\|\xi\|$ passes to limits.
Continuity of multiplication and inversion makes $K$ a compact group.
Average $A^*A$ against its normalized Haar measure. Right invariance of
Haar measure gives (376.1), again with bounds $M^{-2}I,M^2I$.

For finite $G$ this is just
\[
 Q=\frac1{|G|}\sum_{g\in G}\pi(g)^*\pi(g).
\]
The finite-dimensional assertion is therefore no evidence that a group
is amenable. In particular a search through finite matrices alone cannot
produce the required counterrepresentation for a nonamenable group.
Finite matrix calculations are useful for auditing identities, but the
obstruction must occur in an infinite-dimensional representation.

There is a second elementary trap. The matrix
$A=\left(\begin{smallmatrix}1&1\\0&1\end{smallmatrix}\right)$ defines
a representation of $\mathbb Z$ by $n\mapsto A^n$, and every operator
is bounded individually. But $A^n$ has off-diagonal entry $n$, so this
representation is not uniformly bounded. It cannot contradict the
amenable-group theorem or answer the selected problem.

\paragraph{Unitarisability passes to quotients.}
If $G$ is unitarisable and $q:G\twoheadrightarrow H$ is onto, a
uniformly bounded representation $\sigma$ of $H$ pulls back to
$\sigma q$ with the same bound. Any unitariser of the pullback works
for every $\sigma(h)$, because every $h$ has a preimage. Hence $H$
is unitarisable. This elementary fact does not say that unitarisability
passes from a group to an extension of it; the reverse direction is
the one that would be needed to erase the extension in [S1].

\paragraph{Unitarisability passes to arbitrary subgroups.}
\textbf{Lemma 376.2.} If $H\le G$ and $G$ is unitarisable, then $H$
is unitarisable, without restrictions on the index or cardinalities.

\emph{Proof.} Let $\sigma$ be a uniformly bounded representation of
$H$ on $\mathcal K$. Choose representatives $t_x$ for the left cosets
$x\in G/H$, taking $t_H=1$. On
$\mathcal L=\ell^2(G/H;\mathcal K)$ define
\[
 \Pi(g)(\delta_x\xi)
  =\delta_{gx}\,\sigma(t_{gx}^{-1}g t_x)\xi.
 \tag{376.2}
\]
The argument of $\sigma$ belongs to $H$. The identity
\[
 t_{ghx}^{-1}gh t_x
 =(t_{ghx}^{-1}g t_{hx})(t_{hx}^{-1}h t_x)
\]
proves the representation law. Cosets are permuted, and each fiber
operator has norm at most the uniform bound for $\sigma$, so $\Pi$
has that same bound. This computation is valid on finitely supported
vectors and extends by continuity to the Hilbert direct sum, even for
an uncountable index set.

Unitarisability of $G$ supplies an invariant positive $Q$ on
$\mathcal L$ bounded above and below. The fiber at $H$ is invariant
under $\Pi(h)$ for $h\in H$, with action exactly $\sigma(h)$.
Restrict the quadratic form of $Q$ to this fiber. The same two norm
bounds and invariance hold there. The compressed positive operator
therefore satisfies (376.1) for $\sigma$, proving the claim.
\hfill$\square$

It is the restriction of an invariant \emph{form} that is used. An
arbitrary unitarising operator need not preserve this particular fiber,
so simply restricting that operator would not justify the conclusion.

\paragraph{A reduction from arbitrary groups to finitely generated ones.}
We give the corresponding local fact about amenability with its
compactness step visible. If every finitely generated subgroup of $G$
is amenable, then $G$ is amenable.

Consider all means on $\ell^\infty(G)$, a weak-star compact set:
it is the closed set of positive normalized functionals in the dual
unit ball. For any finite list $g_1,\ldots,g_s\in G$, let $H$ be the
subgroup they generate. Take an invariant mean on $H$, and apply it
to restrictions to $H$ of functions on $G$. This gives a mean on $G$
invariant under all $g_i$. Thus the closed conditions of invariance
under individual elements have the finite intersection property.
Compactness supplies a mean invariant under every element of $G$.

Consequently every nonamenable group contains a finitely generated
nonamenable subgroup. If the desired implication were proved for all
finitely generated groups, Lemma 376.2 would extend it to all groups:
every finitely generated subgroup of a unitarisable group would be
amenable, hence the whole group would be amenable. This is a genuine
scope reduction; it does not solve even the finitely generated case.

\paragraph{One representation at a time secretly gives uniform bounds at fixed $M$.}
Let
\[
 c(\pi)=\inf\{\|S\|\|S^{-1}\|:S\text{ unitarises }\pi\}.
\]
If $G$ is unitarisable, then for each fixed $M$ the numbers $c(\pi)$
are bounded over all representations with uniform bound at most $M$.
Suppose otherwise. Choose representations $\pi_j$ of bound at most
$M$ with $c(\pi_j)>j$. Their Hilbert direct sum $\Pi=\bigoplus_j\pi_j$
still has bound at most $M$, so it has an invariant form
$aI\le Q\le bI$. Compressing this form to each invariant summand
gives $c(\pi_j)\le\sqrt{b/a}$ for every $j$, a contradiction.

This does not produce a formula for the bound as a function of $M$,
and it certainly does not reproduce the amenable estimate $M^2$.
The compactness of a single family of forms with \emph{fixed} coercive
bounds can be useful; allowing the lower bounds to tend to zero loses
the invertibility needed in Lemma 376.1. The direct-sum argument identifies
exactly one place where uniformity can be recovered from the global
quantifier in the definition.

\paragraph{A block-matrix route and its exact cohomological obstruction.}
Let $U:G\to\mathcal U(\mathcal H)$ and
$V:G\to\mathcal U(\mathcal K)$ be unitary representations, and let
$b(g):\mathcal K\to\mathcal H$ satisfy
\[
 b(gh)=U(g)b(h)+b(g)V(h).
 \tag{376.3}
\]
Then
\[
 \pi_b(g)=\begin{pmatrix}U(g)&b(g)\\0&V(g)\end{pmatrix}
 \tag{376.4}
\]
is a representation. It is uniformly bounded if and only if the
operators $b(g)$ are uniformly bounded: necessity follows by taking
the off-diagonal compression, and sufficiency follows from the triangle
inequality and the unitary diagonal blocks.

\textbf{Lemma 376.3.} This representation is unitarisable if and only
if some bounded $T:\mathcal K\to\mathcal H$ satisfies
\[
 b(g)=T V(g)-U(g)T\quad(g\in G).
 \tag{376.5}
\]

\emph{Proof.} If (376.5) holds, conjugation by
$S=\left(\begin{smallmatrix}I&-T\\0&I\end{smallmatrix}\right)$
turns (376.4) into $U(g)\oplus V(g)$ by direct multiplication.
Conversely, an invariant equivalent Hilbert inner product makes the
invariant subspace $\mathcal H\oplus0$ have an invariant orthogonal
complement. Orthogonality here is for the new inner product. Since
it is equivalent to the original norm, its projection onto
$\mathcal H\oplus0$ is bounded. The complementary space is therefore
the graph of a bounded operator $T:\mathcal K\to\mathcal H$.
Invariance of that graph says
\[
 \pi_b(g)(T\eta,\eta)=(U(g)T\eta+b(g)\eta,V(g)\eta)
                         =(TV(g)\eta,V(g)\eta),
\]
which is exactly (376.5).\hfill$\square$

For the graph assertion, every $(0,\eta)$ decomposes uniquely as a
vector in $\mathcal H\oplus0$ plus a vector in the complement, so
projection to the second coordinate is a bijection there. Its bounded
inverse follows either from the bounded complement projection or the
bounded inverse theorem. This prevents an unbounded graph operator
from being mistaken for an allowed similarity.

Thus a uniformly bounded solution of (376.3) not of form (376.5) would
give a valid counterrepresentation. Constructing such a solution for
every nonamenable group is the unproved step of this route.

\paragraph{Why a simpler scalar-noise version always collapses.}
For a unitary representation $U$ on a Hilbert space, a bounded cocycle
$c(gh)=c(g)+U(g)c(h)$ gives affine isometries
$A_g(\xi)=U(g)\xi+c(g)$ with a bounded orbit of zero. Every group of
affine isometries of a Hilbert space with a bounded orbit has a fixed
point. Here is a direct proof of the fact used.

Let $B$ be the bounded orbit, and minimize
$F(x)=\sup_{y\in B}\|x-y\|$. A minimizing sequence is bounded. The
parallelogram identity gives
\[
 F((x+z)/2)^2\le\tfrac12F(x)^2+\tfrac12F(z)^2
                         -\tfrac14\|x-z\|^2.
\]
A minimizing sequence is therefore Cauchy and converges to a minimizer;
the same inequality proves uniqueness. Since each affine isometry
permutes $B$, it preserves $F$ and fixes its unique minimizer $x_*$.
Thus $c(g)=x_*-U(g)x_*$ is a coboundary.

This rules out the Hilbert-valued bounded-cocycle shortcut for every
group, even nonamenable ones. The coefficients in (376.3) instead live
in the operator space $\mathcal B(\mathcal K,\mathcal H)$ with its
operator norm, which is not generally a Hilbert space. Applying the
parallelogram argument there without its hypothesis would erase the
very obstruction one is trying to construct.

\paragraph{Exact checks and remaining gap.}
Finite-dimensional rational examples audit (376.1), the averaging
identity and the block similarity formula. The proofs above, rather
than a matrix-size cutoff, establish the general positive results.
No successful finite-dimensional test is advertised as evidence for
the converse, because all such representations are already unitarisable.

\textbf{Outcome: UNRESOLVED.} The amenable implication, subgroup and
quotient permanence, finitely generated reduction and block-matrix
criterion are proved at the stated scopes. The missing step is to use
nonamenability of an arbitrary finitely generated group to create an
infinite-dimensional bounded representation with no coercive invariant
form, for example by an operator-valued cocycle outside (376.5).
The extension theorem in [S1] and the finite-dimensional averaging
argument do not provide that construction for the original group.

\subsection{Sources}
\begin{itemize}
\item[S1]Nicolas Monod and Narutaka Ozawa, The Dixmier problem, lamplighters and Burnside groups, arXiv0902.4585v1 (26Feb2009).
\url{https://arxiv.org/pdf/0902.4585v1}.
\textit{Formulation source.}
\end{itemize}
\end{document}
